\section{Conclusion}
\label{sec:conclusion}

We set out to filter Python SFT training data using doctest execution. We found that only
1 in 1,000 Python files in a curated corpus contains an independently-executable doctest ---
a 31$\times$ gap below the \texttt{>>>} pattern prevalence. Python library code is designed for
installed environments, not clean subprocess execution, and import isolation rejects $\sim$95\%
of AST-parseable doctests at the \texttt{ModuleNotFoundError} stage.

We contribute the first empirical characterization of this gap at corpus scale, validated
with 28/28 unit tests across a three-phase filtering pipeline that processes 10,000 files
in 129.8 seconds. We establish compile-only filtering as the practical execution-quality
gate for raw Python SFT corpus construction and design the SFT quality comparison experiment
as immediate future work.

Measuring feasibility before committing to an experiment design is a small investment that
prevents large experimental waste. The \texttt{>>>} prompt in Python code is documentation
intent, not executable reality --- measuring the gap between the two is the first step toward
principled execution-quality data curation for code LLMs.
